\documentclass{IOS-Book-Article}

\usepackage{mathptmx}
\usepackage{booktabs}
\usepackage{graphicx}
\usepackage{siunitx}
\usepackage{url}
\usepackage{amsmath}

\newcommand{\MatcherLadderCorpusFixes}{42,654}
\newcommand{\MatcherLadderPseudoExamples}{29,411}
\newcommand{\MatcherLadderRows}{%
M1 & Connected baseline & 90.44\% & 32.35\% & 131.8 & 9.225 \\%
M2 & Nearest + street-ref continuity & 93.96\% & 23.22\% & 80.7 & 9.511 \\%
M3 & M2 + connected-candidate gate & 90.76\% & 23.40\% & 131.8 & 9.856 \\%
M4 & Corridor raw mini-HMM & 85.59\% & 17.02\% & 131.8 & 9.521 \\%
M5 & Corridor-aware final & 90.15\% & 53.63\% & 131.8 & 9.093 \\%
M6 & M2 + urban consecutive distance-gap release & 93.96\% & 23.34\% & 80.7 & 8.725 \\%
M7 & M6 + 10m search window & 94.28\% & 25.09\% & 80.7 & 8.488 \\%
M8 & M6 + no-ref street-name continuity & 93.92\% & 23.59\% & 80.7 & 9.429 \\%
M9 & M8 + guarded stale-ref suppression & 93.97\% & 23.59\% & 80.7 & 9.399 \\%
M10 & M9 + node-direction-aware junction release & 93.98\% & 23.65\% & 80.7 & 10.046 \\%
M11 & M10 + topology-free particle sequence & 93.45\% & 25.47\% & 131.8 & 10.087 \\%
M12 & M11 + 10-fix Viterbi sequence & 94.24\% & 26.91\% & 170.9 & 10.374 \\%
Valhalla & Valhalla oracle & 94.04\% & 26.03\% & 155.4 & 7.843 \\
}


\def\hb{\hbox to 11.5 cm{}}

\begin{document}

\hypersetup{pageanchor=false}
\pagestyle{headings}
\def\thepage{}

\begin{frontmatter}

\title{Resource-Constrained On-Device Map Matching for Smartphone Speed-Limit Assistance}

\markboth{}{\hb}

\author[A]{\fnms{Raphael} \snm{Volz}%
\thanks{Corresponding Author: Raphael Volz, Pforzheim University, Germany; E-mail: \textit{raphael.volz@hs-pforzheim.de}.}}

\runningauthor{R. Volz}
\address[A]{Pforzheim University, Germany}

\begin{abstract}
Smartphone speed-limit assistance needs map matching that remains stable under noisy consumer GPS, tunnel loss, and strict on-device artifact budgets. This paper isolates that matcher-design question under a fixed local runtime substrate. A separate companion manuscript submitted to ITSC 2026 studies country-scale runtime data architectures and daily-diff update behavior \cite{volz2026itsc_runtime}; the present paper does not repeat that systems comparison and instead holds the runtime database layer fixed at the currently preferred embedded single-file design. On top of that fixed substrate, we compare a 12-row causal matcher ladder and a desktop Valhalla oracle using closed-loop replay on a 21-log corpus with \MatcherLadderCorpusFixes{} fixes and \MatcherLadderPseudoExamples{} hindsight-labeled examples. On the Karlsruhe seed bundle, the best lightweight row is M7, which reaches 94.28\% replay accuracy and 25.09\% changed-example recall while keeping the 80.7\,MiB no-\texttt{way\_links} bundle tier and reducing query p95 to 8.488\,ms. M12 improves changed-example recall to 26.91\% but requires the 170.9\,MiB full bundle and raises p95 to 10.374\,ms. A Valhalla oracle reaches 94.04\% accuracy and 26.03\% changed-example recall at 7.843\,ms p95 on a 155.4\,MiB tile set. Repeating the replay on a full Baden-W\"urttemberg build leaves the ranking essentially unchanged while increasing native p95 by roughly 3--4\,ms and expanding the Valhalla tiles to 674.5\,MiB. The main result is therefore a deployment recommendation: continue with the lightweight matcher family, specifically the M7 branch, and treat the remaining tunnel and route-stability failures as narrow refinement work rather than a reason to adopt heavier topology tables or a full route engine.
\end{abstract}

\begin{keyword}
intelligent transportation systems\sep intelligent speed assistance\sep map matching\sep mobile GIS\sep offline-first\sep smartphone systems
\end{keyword}
\end{frontmatter}

\markboth{\hb}{\hb}

\section{Introduction}
Intelligent speed assistance is now mandatory in new EU vehicle types, but the installed fleet is far larger than annual new registrations \cite{eu2019_2144,eu2021_1958,eu_mandatory_safety_2024}. That leaves a retrofit space for smartphone-based assistance that can infer speed context from open map data under intermittent connectivity and consumer-grade sensing. In such a product, map matching is necessary but cannot be treated as an unconstrained offline route-inference problem. The runtime must stay auditable, fit inside mobile artifact budgets, and return answers within a single-fix interaction cycle.

This paper studies that deployment-constrained matcher question. It is intentionally narrower than our companion ITSC 2026 submission on runtime data architectures \cite{volz2026itsc_runtime}. That companion paper compares four country-scale storage/query architectures and update behaviors. Here, by contrast, the runtime data layer is held fixed at the preferred embedded single-file design, and only the matching layer is varied. The goal is to identify which causal matcher profile gives the best balance of replay quality, query cost, and bundle size for the current smartphone ISA stack.

The paper makes four contributions:
\begin{enumerate}
\item A fixed 12-row matcher ladder that compares geometric, continuity-aware, topology-bearing, corridor-heavy, and lightweight sequence-aware profiles under one replay protocol.
\item A closed-loop replay benchmark over 21 drive-match logs with \MatcherLadderCorpusFixes{} fixes and \MatcherLadderPseudoExamples{} hindsight-labeled examples.
\item A deployment-facing comparison of lightweight in-house matchers against heavier in-house variants and a desktop Valhalla oracle, scored jointly by replay quality, p95 query time, and bundle size.
\item A seed-versus-full-region sensitivity rerun showing that the model ranking is stable while artifact size and latency penalties scale materially for the heavier options.
\end{enumerate}

\section{Related Work}
Map matching has long been organized as a trade-off between geometric simplicity, topological consistency, and probabilistic sequence reasoning \cite{quddus2007}. HMM-based and path-inference formulations remain the standard answer to noisy or sparse trajectories because later observations can stabilize earlier ambiguous selections \cite{newson2009,lou2009,goh2012,he2013,hunter2014}. Production engines reflect the same principle: OSRM, Mapbox, and Valhalla expose candidate search controls and transition-aware matching interfaces for route-level inference \cite{osrm_match_service,mapbox_map_matching_api,valhalla_meili}.

Smartphone-specific work confirms that these ideas transfer to consumer sensing. Bierlaire et al.\ evaluate a probabilistic matcher directly on smartphone GPS data \cite{bierlaire2013smartphone}. Faster route-level transition checking has also been studied through precomputation-heavy designs such as FMM \cite{yang2018fmm}. Those systems are important references, but they target a larger route-inference budget than the current mobile speed-limit product.

The present work is closer to deployment studies that ask how much matching machinery is actually worth carrying on device. The distinct claim here is not a new HMM or deep model. It is an experimentally grounded recommendation about how far a local smartphone matcher should go before heavier topology, corridor, or route-engine machinery stops paying for itself under product constraints.

\section{Problem Setting}
The deployed YouSpeed runtime performs deterministic way-oriented speed-limit inference against a local OpenStreetMap-derived database. Under the fixed runtime substrate used here, candidate roads are retrieved from a single-file spatially indexed bundle, ranked by local geometry and heading evidence, and then passed through matcher logic that may preserve the current road, promote a same-reference continuation, or release to a new candidate when local evidence is strong enough.

The matcher-design question can be decomposed by which information items each profile is allowed to use:
\begin{itemize}
\item single-fix geometry and heading,
\item continuity state across nearby fixes,
\item same-road or directly connected local topology,
\item short way-sequence or corridor evidence,
\item heavier local graph or route-engine sequence inference.
\end{itemize}

The ladder evaluated in this paper moves through that design space in one direction. M1 is a connected baseline. M2 introduces the lightweight nearest-plus-street-reference continuity rule. M6--M10 refine that lightweight family without reintroducing \texttt{way\_links}. M11 and M12 add stronger local sequence reasoning, and a separate Valhalla oracle provides a reference point for what a production-style route engine can recover when it is allowed to use trace-level matching.

Two constraints define the problem. First, all in-house matching logic must remain compatible with the consumer app's existing way-level speed inference contract. Second, model selection is scored jointly by replay quality, query latency, and artifact size rather than by matching accuracy alone.

\section{Replay Benchmark}
\subsection{Protocol}
The benchmark replays logged fixes in order and rebuilds matcher state from each profile's own previous outputs rather than from the original logged selection. This matters because early mistakes can perturb later state; the protocol therefore evaluates the full recurrent system rather than only pointwise candidate ranking. Hindsight labels are generated by a fixed future-stability rule and applied only where enough later context exists.

The current corpus contains 21 valid drive-match logs with \MatcherLadderCorpusFixes{} replayed fixes and \MatcherLadderPseudoExamples{} hindsight-labeled examples. The reported metrics are overall replay accuracy, changed-example recall, unchanged-example accuracy, p95 query time, and required artifact size.

\subsection{Bundle Tiers and Sensitivity Rerun}
The primary ladder is run on the Karlsruhe seed bundle. Lightweight rows M2 and M6--M10 use an 80.7\,MiB no-\texttt{way\_links} bundle. Topology-bearing rows M1, M3--M5, and M11 use a 131.8\,MiB bundle. M12 requires the 170.9\,MiB full bundle because its sequence path depends on \texttt{way\_endpoints}. The seed Valhalla oracle uses a 155.4\,MiB tile set.

To test whether the seed-region bundle hides a stronger full-region preference, the main competitive rows were rerun on a Baden-W\"urttemberg build. The lightweight in-house base there is a 343.1\,MiB release bundle without the extra topology tables, while the full M12-capable build grows to 733.5\,MiB and the Valhalla tile set to 674.5\,MiB.

\section{Results}
\subsection{Karlsruhe Seed-Bundle Ladder}
Table~1 reports the full Karlsruhe seed-bundle ladder. The central pattern is that heavier machinery does not dominate the lightweight branch once bundle size and p95 latency are included in the decision.

\begin{table*}[!t]
\caption{Closed-loop replay ladder on the Karlsruhe seed bundle (\MatcherLadderCorpusFixes{} fixes; \MatcherLadderPseudoExamples{} hindsight-labeled examples).}
\centering
\scriptsize
\setlength{\tabcolsep}{4pt}
\begin{tabular}{llrrrr}
\toprule
Exp. & Profile & Accuracy & Changed recall & Bundle (MiB) & Query p95 (ms) \\
\midrule
\MatcherLadderRows
\bottomrule
\end{tabular}
\end{table*}

Three rows define the deployment trade-off most clearly. M7 is the best lightweight row at 94.28\% replay accuracy and 25.09\% changed-example recall while keeping the 80.7\,MiB bundle tier and improving p95 to 8.488\,ms. M12 raises changed-example recall to 26.91\%, but the gain is small relative to the added cost: the required artifact grows to 170.9\,MiB and p95 to 10.374\,ms. The Valhalla oracle lands between those two in quality, with 94.04\% accuracy and 26.03\% changed-example recall, but still requires a 155.4\,MiB routing tile set.

The weaker negative results matter as much as the winners. Reintroducing \texttt{way\_links} through M3 does not improve on M2. The corridor-heavy rows M4 and M5 do recover many switch cases, especially M5, but they do so with materially lower overall accuracy. In other words, the broader replay corpus no longer supports the older intuition that moving up the ladder is automatically preferable.

\subsection{Seed Versus Baden-W\"urttemberg}
Table~2 compares the main competitive options between the Karlsruhe seed bundle and the full Baden-W\"urttemberg rerun. The ranking is effectively unchanged: M7 remains the top lightweight operating point, M12 remains the strongest in-house row on changed-example recall, and the Valhalla oracle remains a useful ceiling check rather than a deployment winner.

\begin{table}[!t]
\caption{Seed-versus-Baden-W\"urttemberg sensitivity for the main competitive rows.}
\centering
\scriptsize
\setlength{\tabcolsep}{3pt}
\begin{tabular}{llrrrr}
\toprule
Region & Model & Accuracy & Changed recall & Bundle (MiB) & p95 (ms) \\
\midrule
Seed & M2 & 93.96\% & 23.22\% & 80.7 & 9.511 \\
Seed & M7 & 94.28\% & 25.09\% & 80.7 & 8.488 \\
Seed & M12 & 94.24\% & 26.91\% & 170.9 & 10.374 \\
Seed & Valhalla & 94.04\% & 26.03\% & 155.4 & 7.843 \\
\midrule
BW & M2 & 93.96\% & 23.22\% & 343.1 & 12.927 \\
BW & M7 & 94.28\% & 25.09\% & 343.1 & 12.156 \\
BW & M12 & 94.23\% & 26.91\% & 733.5 & 13.555 \\
BW & Valhalla & 94.06\% & 26.03\% & 674.5 & 7.809 \\
\bottomrule
\end{tabular}
\end{table}

What changes in the full-region rerun is not the ordering but the cost profile. Native in-house p95 rises by roughly 3--4\,ms across the main rows. The heavier artifact classes also widen their size penalty: relative to the Baden-W\"urttemberg lightweight base, the full M12 build is 2.14\(\times\) larger and the Valhalla tiles 1.97\(\times\) larger. Yet these larger artifacts do not produce a commensurate quality gain. The Valhalla oracle changes by only +0.02 percentage points in accuracy between regions, and M12 still improves recall only modestly over M7.

\section{Discussion}
The practical conclusion is narrower than saying that the most complex matcher wins. The present evidence supports three statements.

First, the lightweight branch is no longer just a fallback. M7 strictly dominates M2 on the same bundle tier and also outperforms the topology-bearing and corridor-heavy variants on the current replay corpus. Second, M12 is best understood as a recall-leaning option rather than the default consumer matcher. Its gain on changed-example recall is real, but small relative to the bundle and latency premium. Third, the Valhalla oracle removes the strongest argument for a route-engine pivot. Even with trace-level matching and locate fallback, it does not beat the best in-house row on either deployment-facing target: M7 still leads it on overall accuracy while M12 leads it on changed-example recall.

These results also sharpen the engineering agenda. The remaining open failures are concentrated in tunnel handling and route-stability edge cases rather than in a broad lack of sequence machinery. That means future work should focus on narrow interventions on top of the lightweight branch, for example better tunnel-portal continuity and guarded release rules, rather than reopening the full corridor or route-engine design space.

\section{Conclusion}
This paper carved out the matcher-selection question from the broader YouSpeed system and evaluated it under a fixed mobile runtime substrate. Across a 21-log closed-loop replay benchmark, the best deployment-facing row is M7: it stays on the lightweight no-\texttt{way\_links} bundle tier, improves on M2 in both replay accuracy and changed-example recall, and lowers p95 latency. Heavier sequence-aware variants and a Valhalla oracle recover some additional ambiguous cases, but not enough to justify their larger artifact cost for the current smartphone speed-limit product. The present deployment choice is therefore the lightweight M7 branch, with future work focused on tunnel and route-stability refinements rather than heavier side-table machinery.

\bibliographystyle{vancouver}
\bibliography{icitt2026_refs,../share/references}

\end{document}
